% Generated by scripts/build_reports.py; edit reviewed Markdown inputs.
% Compile with XeLaTeX.
% Options for packages loaded elsewhere
\PassOptionsToPackage{unicode}{hyperref}
\PassOptionsToPackage{hyphens}{url}
\PassOptionsToPackage{dvipsnames,svgnames,x11names}{xcolor}
\documentclass[
  11pt,
]{article}
\usepackage{xcolor}
\usepackage[margin=1in]{geometry}
\usepackage{amsmath,amssymb}
\setcounter{secnumdepth}{-\maxdimen} % remove section numbering
\usepackage{iftex}
\ifPDFTeX
  \usepackage[T1]{fontenc}
  \usepackage[utf8]{inputenc}
  \usepackage{textcomp} % provide euro and other symbols
\else % if luatex or xetex
  \usepackage{unicode-math} % this also loads fontspec
  \defaultfontfeatures{Scale=MatchLowercase}
  \defaultfontfeatures[\rmfamily]{Ligatures=TeX,Scale=1}
\fi
\usepackage{lmodern}
\ifPDFTeX\else
  % xetex/luatex font selection
\fi
% Use upquote if available, for straight quotes in verbatim environments
\IfFileExists{upquote.sty}{\usepackage{upquote}}{}
\IfFileExists{microtype.sty}{% use microtype if available
  \usepackage[]{microtype}
  \UseMicrotypeSet[protrusion]{basicmath} % disable protrusion for tt fonts
}{}
\makeatletter
\@ifundefined{KOMAClassName}{% if non-KOMA class
  \IfFileExists{parskip.sty}{%
    \usepackage{parskip}
  }{% else
    \setlength{\parindent}{0pt}
    \setlength{\parskip}{6pt plus 2pt minus 1pt}}
}{% if KOMA class
  \KOMAoptions{parskip=half}}
\makeatother
\usepackage{longtable,booktabs,array}
\usepackage{caption}
\captionsetup[table]{skip=6pt}
\newcounter{none} % for unnumbered tables
\usepackage{calc} % for calculating minipage widths
% Correct order of tables after \paragraph or \subparagraph
\usepackage{etoolbox}
\makeatletter
\patchcmd\longtable{\par}{\if@noskipsec\mbox{}\fi\par}{}{}
\makeatother
% Allow footnotes in longtable head/foot
\IfFileExists{footnotehyper.sty}{\usepackage{footnotehyper}}{\usepackage{footnote}}
\makesavenoteenv{longtable}
\setlength{\emergencystretch}{3em} % prevent overfull lines
\providecommand{\tightlist}{%
  \setlength{\itemsep}{0pt}\setlength{\parskip}{0pt}}
\usepackage{amsmath,amssymb}
\usepackage{xurl}
\usepackage{fancyhdr}
\pagestyle{fancy}
\fancyhf{}
\fancyhead[L]{\small Pure Thought | Revised research specifications}
\fancyfoot[L]{\small 2026-09-17}
\fancyfoot[R]{\small\thepage}
\setlength{\headheight}{15pt}
\setlength{\emergencystretch}{3em}
\allowdisplaybreaks
\usepackage{bookmark}
\IfFileExists{xurl.sty}{\usepackage{xurl}}{} % add URL line breaks if available
\urlstyle{same}
% fallback for those not using the hyperref driver hyperxmp:
\makeatletter
\@ifundefined{xmpquote}{\newcommand{\xmpquote}[1]{#1}}{}
\makeatother
\hypersetup{
  pdftitle={12. Rigidity and boundary modes of periodic Maxwell frames},
  pdfauthor={Pure Thought research portfolio},
  colorlinks=true,
  linkcolor={blue},
  filecolor={Maroon},
  citecolor={Blue},
  urlcolor={blue},
  pdfcreator={LaTeX via pandoc}}

\title{12. Rigidity and boundary modes of periodic Maxwell frames}
\author{Pure Thought research portfolio}
\date{Revised 17 September 2026}

\begin{document}
\maketitle
\thispagestyle{empty}
\clearpage

\section{12. Rigidity and boundary modes of periodic Maxwell frames}

\textbf{Revised specification --- 2026-09-11.}
\href{https://github.com/rick-stevens-ai/Pure-Thought/blob/main/reviews/12-Topological-Mechanical-Metamaterials.md}{Detailed
critique} \ensuremath{\cdot}
\href{https://github.com/rick-stevens-ai/Pure-Thought/blob/main/archive/original-PRDs/12-Topological-Mechanical-Metamaterials.md}{Original
description} \ensuremath{\cdot}
\href{https://github.com/rick-stevens-ai/Pure-Thought/blob/main/reviews/EVIDENCE-STANDARD.md}{Evidence
standard}

\subsection{Plain-language guide}

\subsubsection{The problem in everyday terms}

A framework of bars or springs can be rigid in one place and easy to
deform in another. Sometimes a carefully arranged repeating framework
has soft motions concentrated at its boundary. The question is whether
we can predict and prove those motions from the exact geometry, rather
than simply observing them in a simulation. Counting connections helps,
but the direction of each connection matters too.

\subsubsection{Key terms}

\begin{itemize}
\item
  \textbf{Frame} --- A network of sites joined by bars or springs, used
  as an ideal mechanical model.
\item
  \textbf{Degree of freedom} --- An independently variable coordinate of
  the system's motion.
\item
  \textbf{Isostatic or Maxwell frame} --- A frame balanced between
  degrees of freedom and constraints in the relevant counting sense;
  details of zero modes and self stress still matter.
\item
  \textbf{Zero mode} --- An infinitesimal displacement requiring no
  restoring energy in the linearized model.
\item
  \textbf{State of self stress} --- A pattern of internal bond tensions
  that balances forces without external loads.
\item
  \textbf{Compatibility matrix C} --- The linear map from small site
  displacements to bond extensions.
\item
  \textbf{Stiffness matrix K} --- The operator relating displacements to
  restoring forces in the linear approximation.
\item
  \textbf{Winding number} --- An integer counting how a complex quantity
  winds around zero along a closed path; here it helps characterize the
  periodic frame.
\item
  \textbf{Termination} --- The precise way a periodic structure is cut
  to make a boundary.
\end{itemize}

\subsubsection{Why this matters}

Geometry-controlled softness could help explain and design structures
with directional compliance, boundary motion or localized mechanical
response. The scientific value is understanding which behaviors follow
from a robust index and which are accidental features of one geometry.
That distinction matters before translating an ideal spring model into a
physical structure.

\subsubsection{What progress would mean}

A useful result would give a frame, an exact count of its allowed soft
motions and a proof of where they concentrate. Material failure,
nonlinear deformation and manufacturing tolerance would remain
additional engineering questions.

\subsection{Core question}

For a geometrically specified isostatic frame and boundary termination,
can the zero-mode count and localization be proved from its
compatibility matrix?

\subsection{Scope and assumptions}

Use unstressed central-force springs with exact rational or algebraic
site coordinates, positive spring constants, and explicit periodic bond
offsets. Fix the unit-cell gauge, dimension, masses, and boundary
termination. Begin with a small frame and a known periodic Maxwell
lattice.

Connectivity alone is insufficient: bond directions determine the
compatibility matrix. Distinguish rigid-body motions, floppy modes, and
states of self stress. Claims concern linearized ideal mechanics.

\subsection{Mathematical target}

Let C map displacements to bond extensions and Q=C\ensuremath{\dagger}.
For a finite free frame,

\[n_0- n_{ss}=dN- N_b,\qquad K=C^{\dagger }\operatorname{diag}(k_b)C,\]

with dynamical matrix M\ensuremath{{}^{-1/2}}KM\ensuremath{{}^{-1/2}}.
For a square periodic C(k), define winding
(2\ensuremath{\pi}i)\ensuremath{{}^{-1}}\ensuremath{\oint}d log det C(k)
only along loops where the determinant is nonzero. Account for acoustic
translation zeros and the local boundary count when translating bulk
winding into an edge index.

\subsection{Required outputs and proof obligations}

Export positions, bonds, C, exact rank/kernel calculations, self-stress
basis, winding certificate, and a boundary-specific mode count. For
localization, provide a transfer-matrix/root-modulus bound separated
from the unit circle. Specify perturbations preserving the nonvanishing
condition and positivity of spring constants.

\subsection{Validation and rejection controls}

Check one elementary underconstrained and one overconstrained frame
against the index equation. Verify K is positive semidefinite and
translations are zero modes where appropriate. Compare two terminations
of the same bulk. Reject winding calculations crossing a zero of det C.

\subsection{Milestones and research extension}

Baseline: exact Maxwell--Calladine counts on finite frames. Strong: one
periodic frame with certified winding and termination-dependent edge
localization. Research extension: a bounded family with robust mode
counts. Elastic beams, prestress, nonlinear stability, and manufacturing
tolerance require separate models.

Milestones are acceptance gates, not calendar promises. Before a run,
freeze the concrete instance/domain, tolerance, compute budget, and
permitted dependencies. A bounded search that fails to find a result
must preserve an unresolved outcome.

\subsection{Starting instruction}

Work on the baseline above. First make every input and convention
explicit, then validate the reference case. Build the checker alongside
the calculation. Label numerical evidence, finite-model certificates,
and continuum theorems separately; return the strongest justified result
and identify the exact unmet proof obligation.

\subsection{Primary references and their role}

\begin{itemize}
\tightlist
\item
  \href{https://arxiv.org/abs/1308.0554}{Kane and Lubensky, Topological
  Boundary Modes in Isostatic Lattices} --- Compatibility matrices,
  topological indices, and boundary modes.
\end{itemize}

References support definitions or limitations; the scope choices and
proposed milestones are editorial recommendations, not claims that these
results have already been obtained.

\subsection{Displayed formulation and claim boundary}

\href{https://github.com/rick-stevens-ai/Pure-Thought/blob/main/PRDs/12-Topological-Mechanical-Metamaterials.md}{Canonical
specification}.

\[Q=C^\dagger,\qquad n_0-n_{ss}=dN-N_b,\qquad
K=C^\dagger\operatorname{diag}(k_b)C,\quad D=M^{-1/2}KM^{-1/2}.\]

\[\nu_\gamma=\frac1{2\pi i}\oint_\gamma d\log\det C(k),
\qquad \det C(k)\ne0\text{ on }\gamma.\]

The first count includes rigid-body motions and self stress. Supply
actual bond directions, unit-cell gauge and termination. Acoustic zeros
and the local boundary count must be included before identifying edge
modes from winding.

\subsection{Detailed review and rationale}

\href{https://github.com/rick-stevens-ai/Pure-Thought/blob/main/PRDs/12-Topological-Mechanical-Metamaterials.md}{Revised
problem} \ensuremath{\cdot}
\href{https://github.com/rick-stevens-ai/Pure-Thought/blob/main/archive/original-PRDs/12-Topological-Mechanical-Metamaterials.md}{Preserved
original} \ensuremath{\cdot}
\href{https://github.com/rick-stevens-ai/Pure-Thought/blob/f854d4e12b55c1f5b5c2a964b75f2d23496d981e/PRDs/12-Topological-Mechanical-Metamaterials.md}{Pinned
upstream source}

Line references below refer to the preserved original at commit
\texttt{f854d4e12b55}, not the rewritten file. ``Blocking'' means the
original issue can invalidate the scientific target or claimed
certificate; ``major'' means it materially changes scope, inference, or
acceptance criteria. These are critiques of specifications, not results
of running the embedded research-code sketches.

\subsubsection{Detailed findings}

\begin{enumerate}
\def\labelenumi{\arabic{enumi}.}
\item
  \textbf{Blocking --- counting formula (22--27, 100).} The original
  replaces the Maxwell--Calladine difference between zero modes and self
  stresses with an incorrect count. It can mislabel even simple rigid
  structures.
\item
  \textbf{Blocking --- mechanical operator (75--78).} The displayed
  dynamical matrix does not provide a reliable central-force
  construction. Building C from bond geometry and
  K=C\ensuremath{\dagger}kC supplies a checkable definition.
\item
  \textbf{Major --- graph-only topology (38, 61--64).} Frames with the
  same graph can have different rigidity and polarization. Geometric
  embedding and unit-cell conventions are part of the input.
\item
  \textbf{Major --- boundary and stability claims (87--103).} A Berry
  phase alone does not give a universal zero-mode count, and a passive
  unstressed stiffness matrix does not have generic
  negative-frequency-squared spectral flow. Boundary indices and
  transfer decay replace the unsupported shortcut.
\end{enumerate}

\subsubsection{What the revision preserves and changes}

The central motivation is retained, but the executable question is now:
\textbf{For a geometrically specified isostatic frame and boundary
termination, can the zero-mode count and localization be proved from its
compatibility matrix?} The revised scope chooses a concrete starting
model; those choices are proposed research-design decisions, not facts
inferred from the original. The baseline can succeed without resolving
the full open research program.

The replacement also removes the original implementation sketches from
the active specification. They remain in the archive for comparison, but
have not been made executable or certified. New implementations should
derive their equations from the revised definitions rather than copy
unchecked prototypes.

\subsubsection{Acceptance criteria to use instead}

Baseline: exact Maxwell--Calladine counts on finite frames. Strong: one
periodic frame with certified winding and termination-dependent edge
localization. Research extension: a bounded family with robust mode
counts. Elastic beams, prestress, nonlinear stability, and manufacturing
tolerance require separate models.

\subsubsection{Sources supporting the corrections}

\begin{itemize}
\tightlist
\item
  \href{https://arxiv.org/abs/1308.0554}{Kane and Lubensky, Topological
  Boundary Modes in Isostatic Lattices} --- Compatibility matrices,
  topological indices, and boundary modes.
\end{itemize}

\section{Evidence and execution standard}

This standard applies to the revised specifications. It is a proposed
contract for future research runs; this editorial revision does not
supply the scientific certificates themselves.

\subsection{What ``pure thought'' means here}

Use explicitly stated mathematical models and deductions without fitting
conclusions to experimental outcomes. Synthetic examples, exact model
inputs, analytic references, numerical search, and established software
are allowed. Physical constants, electronic-structure approximations,
constitutive laws, and empirical parameter sets must be identified when
used. A model-derived result is conditional on that model.

The original ban on legacy software conflicts with its proposed reliance
on NumPy, SciPy, CVXPY, nauty, and other libraries. Reimplementing every
optimizer and algebra system is not necessary for scientific
independence. Prefer a small independently checkable certificate and a
documented trusted base. A fresh implementation can be a project
objective, but must not be confused with a mathematical proof or novelty
claim.

\subsection{Four evidence levels}

{\def\LTcaptype{none} % do not increment counter
\begin{longtable}[]{@{}
  >{\raggedright\arraybackslash}p{(\linewidth - 4\tabcolsep) * \real{0.3333}}
  >{\raggedright\arraybackslash}p{(\linewidth - 4\tabcolsep) * \real{0.3333}}
  >{\raggedright\arraybackslash}p{(\linewidth - 4\tabcolsep) * \real{0.3333}}@{}}
\toprule\noalign{}
\begin{minipage}[b]{\linewidth}\raggedright
Label
\end{minipage} & \begin{minipage}[b]{\linewidth}\raggedright
What it establishes
\end{minipage} & \begin{minipage}[b]{\linewidth}\raggedright
What it does not establish
\end{minipage} \\
\midrule\noalign{}
\endhead
\bottomrule\noalign{}
\endlastfoot
Numerical evidence & A reproducible computation at stated precision,
sample size, and discretization & Exact truth, continuum coverage, or a
universal theorem \\
Finite-model certificate & An exact/validated statement about the
specified finite model or finite relaxation & An omitted tail,
thermodynamic limit, or physical realization \\
Conditional mathematical theorem & A proof over the stated domain with
every hypothesis and error controlled & Truth outside those assumptions
or a demonstrated physical system \\
Formal proof & A checked theorem in a named proof assistant, with
imported axioms/dependencies listed & Correct physical modeling merely
by virtue of formalization \\
\end{longtable}
}

Formalization is an optional extra deliverable unless the specific
project makes it mandatory. High precision, JSON output, solver status
and small residuals do not by themselves confer any of these proof
levels.

\subsection{Claim directions}

For minimization over a larger feasible set, the optimum is a lower
bound on the original minimum. A physical feasible point supplies an
upper bound. For maximization the directions reverse. Track feasibility
of the primal/dual witnesses separately from their objective values.

Infeasibility of an inner ansatz excludes only that ansatz.
Infeasibility of an outer relaxation can exclude the original problem if
the map into the relaxation is proved and all approximations preserve
necessity. Passing necessary tests means ``not excluded''; existence
needs a constructive sufficient argument.

A smallest observed example is not a minimality theorem. A complete
catalog must define the universe, equivalence relation, boundaries, and
proof that pruning misses nothing. A finite-size numerical trend does
not prove an asymptotic limit.

\subsection{Minimum run manifest}

Before calculation, record the problem ID and revision hash, exact
inputs and conventions, quantified domain, equivalence relation, target
observable, accepted evidence level, tolerances, and resource limits.
Record dependencies and precision/rounding mode. Where randomization is
used, preserve seeds, trial counts and raw sufficient statistics.
Resource limits are planning decisions, not evidence that a mathematical
solution cannot exist.

An output bundle should include:

\begin{itemize}
\tightlist
\item
  The statement being claimed, including all assumptions and
  quantifiers.
\item
  Construction or proof object, plus an independently runnable verifier
  and its trust assumptions.
\item
  Numerical/error ledger separating roundoff, discretization,
  truncation, optimization, sampling and modeling error where relevant.
\item
  Positive reference cases and negative controls that the verifier must
  reject.
\item
  Coverage: included domain, omitted domain, unresolved regions and
  termination reason.
\end{itemize}

These are requirements for future implementations, not a promise of
polynomial-size certificates for every problem.

\subsection{Checking specific certificate types}

\textbf{LP/SDP/SOS:} verify equality constraints, objective direction,
and cone feasibility using exact arithmetic or outward-rounded
intervals. Rationalizing a numerical solution can break PSD constraints.
Correct residuals with justified norm bounds; do not silently project
and reuse the old objective. Distinguish a proof of a polynomial
approximation from a proof for the original function.

\textbf{Interval numerics:} state function enclosures, domains and
outward rounding. A small root box needs existence and, if claimed,
uniqueness. Cover the complement for a global enumeration. Include
singularities and tails rather than discarding them numerically.

\textbf{SAT/SMT:} a solver's UNSAT label is not a certificate. Preserve
a proof supported by a checker and the proof that the encoding, symmetry
reduction and blocked solutions match the mathematical universe. DRAT
applies to an appropriate Boolean CNF proof, not automatically to every
real-arithmetic or SDP claim.

\textbf{Statistics:} specify the estimand, dependence assumptions,
confidence procedure and bias. Zero observed failures does not prove
zero failure probability. MCMC error bars based on estimated
autocorrelation are not generally rigorous mixing guarantees.

\subsection{Milestones and novelty}

Baseline means a working, independently validated result on a concrete
reference problem. Strong means a broader or sharper result with its
claimed evidence level. Research extension means a possible novel result
whose feasibility and significance need review. A negative theorem can
be valuable; an unsuccessful search can also be useful if its scope is
documented, but the two are not interchangeable.

Do not promise ``publication-quality'' in a fixed number of months.
Before claiming novelty, conduct a focused current literature comparison
against the exact theorem/model/metric. The sources in this revision
check key definitions and limitations; they are not an exhaustive
novelty audit across all 30 fields.

\end{document}
